% ============================================================
\section{Motivations}\label{sec:motivations}
% ============================================================
Le problème de la recherche d'une partition satisfaisante trouve des échos directs dans l'analyse des structures relationnelles modernes. 

En \textbf{recherche de communautés} (réseaux sociaux ou biologiques), les sommets représentent des individus ou des entités biologiques (gènes, protéines) et les arêtes correspondent à leurs interactions. 
Une partition satisfaisante isole des sous-groupes stables où l'activité endogène prédomine sur les échanges exogènes.
Cela permet d'identifier des communautés d'intérêt, comme des groupes d'amis dans un réseau social ou des groupes de gènes fonctionnellement liés dans un réseau biologique.

En \textbf{théorie des jeux et sociologie}, ce problème modélise la formation d'alliances défensives. 
Un groupe de sommets \(A\) forme une alliance stable si chaque membre trouve au moins autant de soutien au sein de \(A\) qu'il n'affronte de menaces à l'extérieur (\(B\)). 
Cette modélisation s'applique historiquement aux équilibres géopolitiques et industriels (fusions et alliances stratégiques).

Enfin, en \textbf{intelligence artificielle}, le partitionnement de graphes est un pilier de l'apprentissage non supervisé (\textit{clustering}). 
Garantir qu'un cluster possède une densité interne supérieure à sa connectivité externe améliore la robustesse des algorithmes de classification textuelle et de vision par ordinateur.

L'étude de ce problème, à la fois théorique et appliquée, est donc cruciale pour comprendre les liens complexes dans les réseaux modernes et pour développer des outils d'analyse dans des domaines aussi variés que la sociologie, la biologie et l'informatique.